\section{Related Work}

\paragraph{Transformer forecasting backbones.}
Transformer forecasting models adapt sequence attention to long input windows and multivariate horizons. Informer \citep{zhou2021informer} uses ProbSparse attention and an encoder-decoder design to reduce the cost of long-sequence forecasting. PatchTST \citep{nie2023patchtst} changes the representation by treating a time series as a sequence of patches, while iTransformer \citep{liu2023itransformer} inverts the tokenization pattern so variables rather than time steps become the tokens. These works mainly change the architecture or representation. Our study is orthogonal: we keep the backbone fixed within each experiment and isolate the activation family and placement.

\paragraph{Activation functions as controlled design variables.}
ReLU \citep{nair2010relu} and GELU \citep{hendrycks2016gelu} are common neural-network defaults. GELU is particularly common in transformer implementations because it is smooth and has become part of many standard feed-forward blocks. In long-horizon forecasting code, however, the default activation is rarely treated as an experimental factor on the same footing as attention, width, patch length, or decomposition. Our results show why this is risky: under a matched Informer protocol, replacing the activation alone changes which configuration is best in nearly every dataset-horizon cell.

\paragraph{Bounded and periodic nonlinearities.}
Bounded nonlinearities such as \texttt{tanh} and \texttt{softsign} constrain activation magnitude while preserving sign. Periodic activations and Fourier-style features have also been used to represent high-frequency functions and signals \citep{sitzmann2020siren,tancik2020fourier}. Our oscillatory \texttt{tanh\_sin001} variant is much more modest than the periodic activations used in implicit neural representations: it adds a small sinusoidal perturbation around \texttt{tanh} inside a forecasting transformer's feed-forward blocks. This makes the comparison against the unperturbed bounded baseline essential. The full grid shows that \texttt{tanh\_sin001} is a strong static candidate, but it does not justify a general claim that oscillation is always better than \texttt{tanh}.

\paragraph{Regime-aware forecasting diagnostics.}
Forecasting errors depend on input-window regimes such as volatility, jumps, trend consistency, and seasonal concentration. It is therefore tempting to turn activation selection into a zero-shot routing problem: compute factors from the input window and choose an activation or placement. Our current evidence is more conservative. Factor bins reveal that activation differences vary across regimes, but retrospective bin contrasts are not deployment evidence. We therefore distinguish three levels of evidence: full-grid static performance, same-split diagnostic oracles, and leave-one-seed-out or held-out router validation. Only the first is currently strong enough to support the main paper claim.

\paragraph{Architecture transfer.}
If activation sensitivity were purely a property of the dataset, an activation ranking found in Informer should transfer directly to PatchTST or iTransformer. The transfer tests contradict that simple story. PatchTST shows selected non-GELU wins but keeps GELU as the strongest static choice in most tested cells. iTransformer shows even weaker FFN-only transfer in the current one-seed pilot. These negative and partial results are important because they define the scope of the contribution: activation choice matters, but it is jointly conditioned on dataset, horizon, placement, and architecture.
